\section{Correctness and survival of the nested sets}\label{sec:correctness}

We first analyze the estimation-free recursion.  For integers $d,s\ge0$, let
$P_d(s)$ be the infimum, over all graphs, candidate sets of size $s$, and
distinguished pivots in those sets, of the probability that the ideal
size-biased recursion completes $d$ additional steps.  Set $P_0(s)=1$ for
$s\ge1$, $P_d(0)=0$, and $P_d(1)=0$ for $d\ge1$.

\begin{lemma}[Size-biased survival]\label{lem:size-biased-survival}
For every $d\ge0$ and $s\ge1$,
\[
 P_d(s)\ge \frac{(s-2^d+1)_+}{s},
 \qquad (y)_+=\max\{y,0\}.
\]
In particular, for $M=4^{K-1}$ and $r=2K-3$, the ideal recursion completes
with probability at least $1/2+1/M$.
\end{lemma}

\begin{proof}
The case $d=0$ is immediate.  Suppose $d\ge1$ and $s\ge2$, and let the two
colour classes among the $s-1$ non-pivot vertices have sizes $a$ and $b$.
A uniform sample selects these classes with probabilities $a/(s-1)$ and
$b/(s-1)$.  The induction hypothesis therefore gives
\[
 P_d(s)\ge
 \frac{aP_{d-1}(a)+bP_{d-1}(b)}{s-1}
 \ge
 \frac{(a-2^{d-1}+1)_+ +(b-2^{d-1}+1)_+}{s-1}.
\]
Because $(x)_++(y)_+\ge(x+y)_+$ and $a+b=s-1$, the numerator is at least
$(s-2^d+1)_+$.  Replacing the denominator $s-1$ by $s$ weakens the bound and
proves the claim; the definitions cover $s\le1$.  Finally, $M=2^{r+1}$, so
the bound at $(d,s)=(r,M)$ is
$(M-2^r+1)/M=1/2+1/M$.
\end{proof}

Exact conditional uniformity of each successful capped search couples an
implemented run to the ideal recursion.  Assign each of the $r$ searches
failure probability at most $1/(100r)$.  A union bound decreases the
completion probability by at most $1/100$, so a single fixed-cap run succeeds
with probability at least $0.49$.  Independent repetition, followed by exact
verification, amplifies this probability to $1-\eta$ using
$O(\log(1/\eta))$ runs.  On every completed run, nesting gives
$G(v_i,v_j)=G(v_i,w)=c_i$ for $j>i$; the pigeonhole argument in
\cref{lem:output} then gives a homogeneous output.

We next analyze the scale-aware recursion.  The first lemma controls all
adaptive estimates with one event.

\begin{lemma}[Conditional near-majorities]\label{lem:concentration}
Except with probability at most $\eta/2$, every completed sampling batch
reached after successful searches and an accurate prefix satisfies
\[
 |\widehat p_i-p_i|\le\varepsilon_i,
\]
where $p_i$ is the true colour-$1$ fraction from $v_i$ to
$S_i\setminus\{v_i\}$.
\end{lemma}

\begin{proof}
Condition on the complete classical history $\mathcal H_i$ before the samples
at level $i$.  The graph, $S_i$, and $v_i$ are then fixed.  Let
$\mathcal B_i$ be the event that the verified batch succeeds.  Conditional on
$\mathcal B_i$, \cref{lem:capped-search} gives independent uniform samples
from the fixed set.  Hoeffding's inequality \cite{Hoeffding1963} and
\cref{eq:sample-count} imply
\[
 \Pr\!\left(
 |\widehat p_i-p_i|>\varepsilon_i\mid\mathcal H_i,\mathcal B_i
 \right)
 \le2e^{-2m_i\varepsilon_i^2}
 \le\frac{\eta}{2r}.
\]
Multiplying by $\Pr(\mathcal B_i\mid\mathcal H_i)\le1$ gives the same upper
bound for a completed but inaccurate batch conditional on $\mathcal H_i$.
Apply this bound to the first inaccurate completed level after an accurate
prefix.  The tower property removes the conditioning, and a union bound over
the $r$ possible levels completes the proof.  The argument does not require
independence across levels.
\end{proof}

There are $r$ pivot searches, $r$ sampling batches, and one final search.
Give each operation failure probability at most $\eta/(4r+2)$ using the
capped amplification in \cref{lem:capped-search}.  Applying these conditional
bounds after good prefixes, as above, shows that their union has probability
at most $\eta/2$.  Together with \cref{lem:concentration}, all searches and
estimates used below are good with probability at least $1-\eta$.

Let $s_i=|S_i|$ and
\[
 A_0=1,
 \qquad A_i=\prod_{j=0}^{i-1}a_j.
\]

\begin{lemma}[Variable recurrence and density]\label{lem:survival}
On every prefix satisfying \cref{lem:concentration},
\[
 s_i>MA_i-1,
 \qquad
 A_i\ge\frac78\,2^{-i}.
\]
Consequently $S_r$ is nonempty, and for every sampling level $i<r$,
\[
 \frac{|S_i\setminus\{v_i\}|}{M}
 >\frac38\,2^{-i}.
\]
The final set also satisfies
$|S_r|/M\ge2^{-(r+1)}\ge(3/8)2^{-r}$.
\end{lemma}

\begin{proof}
An accurate empirical majority retains at least the fraction
$a_i=1/2-\varepsilon_i$ of the non-pivot vertices.  Hence
\begin{equation}\label{eq:variable-recurrence}
 s_{i+1}\ge a_i(s_i-1).
\end{equation}
Since $\sum_i\varepsilon_i=1/16$ and
$\prod_j(1-x_j)\ge1-\sum_jx_j$ for $x_j\in[0,1]$,
\[
 A_i
 =2^{-i}\prod_{j<i}(1-2\varepsilon_j)
 \ge2^{-i}\left(1-2\sum_{j<i}\varepsilon_j\right)
 \ge\frac78\,2^{-i}.
\]

We prove $s_i>MA_i-1$ by induction.  The claim holds at $i=0$.  If it holds
at level $i$, then \cref{eq:variable-recurrence} and $a_i<1/2$ give
\[
 s_{i+1}>a_i(MA_i-2)=MA_{i+1}-2a_i>MA_{i+1}-1.
\]
Because $M2^{-r}=2$, we obtain
$s_r>(7/8)2-1=3/4$, and therefore the integer $s_r$ is positive.  For $i<r$,
we have $M2^{-i}\ge4$ and
\[
 s_i-1>MA_i-2
 \ge\frac78M2^{-i}-2
 \ge\frac38M2^{-i}.
\]
Finally, $s_r\ge1$ and $1/M=2^{-(r+1)}$, which proves the last assertion.
\end{proof}

\begin{lemma}[Exact homogeneous output]\label{lem:output}
Whenever the algorithm reaches its final verified output, that output is a
$K$-clique or a $K$-vertex independent set.
\end{lemma}

\begin{proof}
For $i<j$, the definition of $S_{i+1}$ gives
$G(v_i,v_j)=c_i$, and it also gives $G(v_i,w)=c_i$.  Among the
$r=2K-3$ colours, one occurs at least
$\lceil r/2\rceil=K-1$ times.  The corresponding pivots together with $w$
therefore induce that colour.  The final direct check verifies the returned
set independently of earlier search or estimation errors.
\end{proof}
